\documentclass{article}
\textheight 23.5cm \textwidth 15.8cm
%\leftskip -1cm
\topmargin -1.5cm \oddsidemargin 0.3cm \evensidemargin -0.3cm
%\documentclass[final]{siamltex}

\usepackage{graphicx}
\usepackage{epsfig}
\usepackage{amsmath}
\usepackage{amsthm}
\usepackage{amssymb}
\usepackage{mathrsfs}
\usepackage{float}
\usepackage{multirow}
\usepackage{verbatim}
\usepackage{fancyhdr}
\usepackage{subfigure}
\usepackage{color}
\usepackage{mathtools}
\usepackage{mathrsfs}
%\usepackage{natbib}
\usepackage{sectsty}
%\usepackage[title]{appendix}
\usepackage{threeparttable}
\usepackage{dcolumn}
\usepackage{booktabs}
\usepackage{indentfirst}
\usepackage{setspace}
\usepackage{bm}
\usepackage{enumerate}
\usepackage{geometry}

\title{Problem 1.9: Inexact Line Search for the Rosenbrock Function}
\author{Tianyang Li}

\begin{document}
\maketitle

\section{Introduction}

This problem considers the inexact one-dimensional line search
\[
    \min_{\alpha \ge 0} \varphi(\alpha) = f(\bm{x}_k + \alpha \bm{p}_k)
\]
for the Rosenbrock function
\[
    f(x_1, x_2) = 100(x_2 - x_1^2)^2 + (1 - x_1)^2.
\]
The given point and search direction are
\[
    \bm{x}_k = (-1, 1)^T, \qquad \bm{p}_k = (1, 1)^T.
\]
The goal is to determine a suitable step size $\alpha$ along this direction.

\section{Method}

Substituting
\[
    \bm{x}_k + \alpha \bm{p}_k = (-1 + \alpha, 1 + \alpha)^T
\]
into the Rosenbrock function gives the one-dimensional objective
\[
    \varphi(\alpha) = 100\alpha^4 - 600\alpha^3 + 901\alpha^2 - 4\alpha + 4.
\]
Its derivative is
\[
    \varphi'(\alpha) = 400\alpha^3 - 1800\alpha^2 + 1802\alpha - 4.
\]

At the starting point,
\[
    \nabla f(-1, 1) = (-4, 0)^T,
\]
so the directional derivative is
\[
    \nabla f(\bm{x}_k)^T \bm{p}_k = (-4, 0)
    \begin{pmatrix}
        1 \\
        1
    \end{pmatrix}
    = -4 < 0.
\]
Therefore $\bm{p}_k$ is a descent direction.

For the inexact line search, I use the Armijo backtracking rule with
\[
    \alpha_0 = 1, \qquad \rho = 0.5, \qquad c_1 = 10^{-4}.
\]
The Armijo condition is
\[
    \varphi(\alpha) \le \varphi(0) + c_1 \alpha \varphi'(0).
\]
Since $\varphi(0) = 4$ and $\varphi'(0) = -4$, this becomes
\[
    \varphi(\alpha) \le 4 - 0.0004\alpha.
\]

\section{Result}

As a reference, solving $\varphi'(\alpha) = 0$ on $\alpha \ge 0$ gives three stationary points:
\[
    \alpha \approx 0.002225, \qquad \alpha \approx 1.498886, \qquad \alpha \approx 2.998889.
\]
Their corresponding function values are
\[
    \varphi(0.002225) \approx 3.995554,
\]
\[
    \varphi(1.498886) \approx 506.500557,
\]
\[
    \varphi(2.998889) \approx 0.998889.
\]
Hence the exact line-search minimizer along this ray is
\[
    \alpha_* \approx 2.998889.
\]

Using Armijo backtracking, the tested step sizes are shown in Table~\ref{tab:armijo}.

\begin{table}[H]
    \centering
    \begin{tabular}{ccccc}
        \toprule
        Iteration & $\alpha$ & $\varphi(\alpha)$ & Armijo right-hand side & Accept \\
        \midrule
        1         & 1.000000 & 401.000000        & 3.999600               & no     \\
        2         & 0.500000 & 158.500000        & 3.999800               & no     \\
        3         & 0.250000 & 50.328125         & 3.999900               & no     \\
        4         & 0.125000 & 16.430664         & 3.999950               & no     \\
        5         & 0.062500 & 7.124573          & 3.999975               & no     \\
        6         & 0.031250 & 4.736668          & 3.999988               & no     \\
        7         & 0.015625 & 4.155188          & 3.999994               & no     \\
        8         & 0.007812 & 4.023457          & 3.999997               & no     \\
        9         & 0.003906 & 3.998087          & 3.999998               & yes    \\
        \bottomrule
    \end{tabular}
    \caption{Armijo backtracking line search results.}
    \label{tab:armijo}
\end{table}

Therefore the accepted inexact step size is
\[
    \alpha = 0.003906.
\]
The new point is
\[
    \bm{x}_{k+1} = \bm{x}_k + \alpha \bm{p}_k
    = (-1, 1)^T + 0.003906(1, 1)^T
    \approx (-0.996094, 1.003906)^T.
\]
At this point,
\[
    f(\bm{x}_{k+1}) = \varphi(0.003906) \approx 3.998087.
\]

\section{Conclusion}

For the Rosenbrock function at $\bm{x}_k = (-1,1)^T$ along the direction $\bm{p}_k = (1,1)^T$, the induced one-dimensional objective is a quartic polynomial. The direction is indeed a descent direction because the directional derivative is negative at the starting point. Using Armijo backtracking with $\alpha_0 = 1$, $\rho = 0.5$, and $c_1 = 10^{-4}$, the first acceptable inexact line-search step is
\[
    \alpha = 0.003906.
\]
This step decreases the objective from $4$ to approximately $3.998087$. For comparison, the exact minimizer along the same ray is much larger, namely $\alpha_* \approx 2.998889$, with a smaller function value of about $0.998889$. This shows that inexact line search produces a conservative but valid descent step without solving the one-dimensional minimization problem exactly.

\end{document}